% Part: first-order-logic
% Chapter: tableaux
% Section: proving-things-quant

\documentclass[../../../include/open-logic-section]{subfiles}

\begin{document}

\olfileid{fol}{tab}{prq}

\olsection{પરિમાણકો સાથેના \usetoken{P}{tableau}}

\begin{ex}
પરિમાણકો સાથે કામ કરતી વખતે આઇગનચલ-શરતનો ભંગ ન થાય તેની ખાતરી કરવી
પડે છે અને તે માટે કેટલીક વાર અમુક અનુમાનો કરવાના ક્રમમાં ફેરફાર કરવો
પડે છે. સામાન્ય રીતે આઇગનચલ-શરતને આધીન નિયમોને પહેલાં સંભાળવાનો પ્રયત્ન
કરવાથી મદદ મળે છે (તેઓ પૂર્ણ થયેલા !!{tableau}માં વધુ ઉપર આવશે).

નીચેના !!{sentence} માટે આપણે !!a{tableau} કેવી રીતે આપીશું તે જોઈએ:
$\lexists[x][\lnot !A(x)] \lif \lnot \lforall[x][!A(x)]$.
હંમેશની જેમ, ધારણા નોંધીને શરૂઆત કરીએ છીએ:
\begin{oltableau}
[\sFmla{\False}{\lexists[x][\lnot \formula{A}(x)]\lif \lnot
    \lforall[x][\formula{A}(x)]}, just=\TAss]
\end{oltableau}
!!{main operator} $\lif$ હોવાથી આપણે $\TRule{\False}{\lif}$ લાગુ કરીએ
છીએ:
\begin{oltableau}
[\sFmla{\False}{\lexists[x][\lnot \formula{A}(x)]\lif \lnot
    \lforall[x][\formula{A}(x)]}, just=\TAss, checked
  [\sFmla{\True}{\lexists[x][\lnot \formula{A}(x)]},
    just={\TRule{\False}{\lif}[1]}
    [\sFmla{\False}{\lnot\lforall[x][\formula{A}(x)]},
      just={\TRule{\False}{\lif}[1]}]
  ]
]
\end{oltableau}
હવે સંભાળવાની પંક્તિ~$2$ છે. આપણે $\TRule{\True}{\lexists}$ વાપરીએ
છીએ. તેના માટે નવો !!{constant} જોઈએ; હજી કોઈ !!{constant}s આવતાં ન
હોવાથી ગમે તે પસંદ કરી શકીએ છીએ, કહો કે~$a$.
\begin{oltableau}
[\sFmla{\False}{\lexists[x][\lnot \formula{A}(x)]\lif \lnot
    \lforall[x][\formula{A}(x)]}, just=\TAss, checked
  [\sFmla{\True}{\lexists[x][\lnot \formula{A}(x)]},
    just={\TRule{\False}{\lif}[1]}, checked
    [\sFmla{\False}{\lnot\lforall[x][\formula{A}(x)]},
      just={\TRule{\False}{\lif}[1]}
      [\sFmla{\True}{\lnot\formula{A}(a)}, just={\TRule{\True}{\lexists}[2]}]
    ]
  ]
]
\end{oltableau}
હવે પંક્તિ~$3$ને $\TRule{\False}{\lnot}$ લાગુ કરીએ:
\begin{oltableau}
[\sFmla{\False}{\lexists[x][\lnot \formula{A}(x)]\lif \lnot
    \lforall[x][\formula{A}(x)]}, just=\TAss, checked
  [\sFmla{\True}{\lexists[x][\lnot \formula{A}(x)]},
    just={\TRule{\False}{\lif}[1]}, checked
    [\sFmla{\False}{\lnot\lforall[x][\formula{A}(x)]},
      just={\TRule{\False}{\lif}[1]}, checked
      [\sFmla{\True}{\lnot\formula{A}(a)}, just={\TRule{\True}{\lexists}[2]}
        [\sFmla{\True}{\lforall[x][\formula{A}(x)]},
          just = {\TRule{\False}{\lnot}[3]}
        ]
      ]
    ]
  ]
]
\end{oltableau}
પંક્તિ~$4$ને $\TRule{\True}{\lnot}$ અને ત્યાર પછી પંક્તિ~$5$ને
$\TRule{\True}{\lforall}$ લાગુ કરતાં બંધ !!{tableau} મળે છે.
\begin{oltableau}
[\sFmla{\False}{\lexists[x][\lnot \formula{A}(x)]\lif \lnot
    \lforall[x][\formula{A}(x)]}, just=\TAss, checked
  [\sFmla{\True}{\lexists[x][\lnot \formula{A}(x)]},
    just={\TRule{\False}{\lif}[1]}, checked
    [\sFmla{\False}{\lnot\lforall[x][\formula{A}(x)]},
      just={\TRule{\False}{\lif}[1]}, checked
      [\sFmla{\True}{\lnot\formula{A}(a)}, just={\TRule{\True}{\lexists}[2]}
        [\sFmla{\True}{\lforall[x][\formula{A}(x)]},
          just = {\TRule{\False}{\lnot}[3]}
          [\sFmla{\False}{\formula{A}(a)}, just={\TRule{\True}{\lnot}[4]}
            [\sFmla{\True}{\formula{A}(a)},
              just={\TRule{\True}{\lforall}[5]}, close
            ]
          ]
        ]
      ]
    ]
  ]
]
\end{oltableau}
\end{ex}

\begin{ex}
નીચેના ગણ માટે આપણે !!a{tableau} કેવી રીતે આપીશું તે જોઈએ:
\[
\sFmla{\False}{\lexists[x][!C(x,b)]},
\sFmla{\True}{\lexists[x][(!A(x) \land !B(x))]},
\sFmla{\True}{\lforall[x][(!B(x) \lif !C(x,b))]}.
\]
હંમેશની જેમ, ધારણાઓથી શરૂઆત કરીએ છીએ:
\begin{oltableau}
  [\sFmla{\False}{\lexists[x][\formula{C}(x,b)]}, just=\TAss
    [\sFmla{\True}{\lexists[x][(\formula{A}(x) \land \formula{B}(x))]},
      just=\TAss
      [\sFmla{\True}{\lforall[x][(\formula{B}(x) \lif \formula{C}(x,b))]},
        just=\TAss]
    ]
  ]
\end{oltableau}
આપણે આઇગનચલ-શરત ધરાવતો નિયમ હંમેશાં પહેલાં લાગુ કરવો જોઈએ; આ કિસ્સામાં
તે પંક્તિ~$2$ને $\TRule{\True}{\lexists}$ લાગુ કરવાનું છે. ધારણાઓમાં
!!{constant}~$b$ આવતો હોવાથી જુદો અચળ વાપરવો પડે છે; ફરીથી~$a$ પસંદ
કરીએ.
\begin{oltableau}
  [\sFmla{\False}{\lexists[x][\formula{C}(x,b)]}, just=\TAss
    [\sFmla{\True}{\lexists[x][(\formula{A}(x) \land \formula{B}(x))]},
      just=\TAss, checked
      [\sFmla{\True}{\lforall[x][(\formula{B}(x) \lif \formula{C}(x,b))]},
        just=\TAss
        [\sFmla{\True}{\formula{A}(a) \land \formula{B}(a)},
          just={\TRule{\True}{\lexists}[2]}]
      ]
    ]
  ]
\end{oltableau}
હવે પંક્તિ~$1$ને $\TRule{\False}{\lexists}$ અથવા પંક્તિ~$3$ને
$\TRule{\True}{\lforall}$ લાગુ કરીએ તો $x$ની જગ્યાએ કયું પદ~$t$ મૂકવું
તે નક્કી કરવું પડે છે. આ નિયમોને આઇગનચલ-શરત ન હોવાથી ગમે તે પદ પસંદ
કરી શકીએ છીએ. કેટલાક કિસ્સામાં જુદા જુદા~$t$ લઈને નિયમ ઘણી વાર પણ
લાગુ કરવો પડે. પરંતુ સામાન્ય નિયમ તરીકે !!{tableau}માં પહેલેથી આવતાં
પદોમાંથી કોઈ એક પસંદ કરવું લાભદાયી છે---અહીં $a$ અને~$b$ આવે છે---અને
આ કિસ્સામાં આપણે અનુમાન કરી શકીએ કે $a$થી બંધ શાખા મળવાની સંભાવના વધુ
છે.
\begin{oltableau}
  [\sFmla{\False}{\lexists[x][\formula{C}(x,b)]}, just=\TAss
    [\sFmla{\True}{\lexists[x][(\formula{A}(x) \land \formula{B}(x))]},
      just=\TAss, checked
      [\sFmla{\True}{\lforall[x][(\formula{B}(x) \lif \formula{C}(x,b))]}, just=\TAss
        [\sFmla{\True}{\formula{A}(a) \land \formula{B}(a)}, just={\TRule{\True}{\lexists}[2]}
          [\sFmla{\False}{\formula{C}(a,b)}, just={\TRule{\False}{\lexists}[1]}
            [\sFmla{\True}{\formula{B}(a) \lif \formula{C}(a,b)},
              just={\TRule{\True}{\lforall}[3]}
            ]
          ]
        ]
      ]
    ]
  ]
\end{oltableau}
પંક્તિઓ $1$ અને $3$નાં !!{signed formula}s ફરી વાપરવાં પડે તેમ હોવાથી
આપણે તેમની બાજુમાં ખરાનું ચિહ્ન મૂકતા નથી. હવે પંક્તિ~$4$ને
$\TRule{\True}{\land}$ લાગુ કરો:
\begin{oltableau}
  [\sFmla{\False}{\lexists[x][\formula{C}(x,b)]}, just=\TAss
    [\sFmla{\True}{\lexists[x][(\formula{A}(x) \land \formula{B}(x))]},
      just=\TAss, checked
      [\sFmla{\True}{\lforall[x][(\formula{B}(x) \lif \formula{C}(x,b))]}, just=\TAss
        [\sFmla{\True}{\formula{A}(a) \land \formula{B}(a)},
          just={\TRule{\True}{\lexists}[2]}, checked
          [\sFmla{\False}{\formula{C}(a,b)}, just={\TRule{\False}{\lexists}[1]}
            [\sFmla{\True}{\formula{B}(a) \lif \formula{C}(a,b)},
              just={\TRule{\True}{\lforall}[3]}
              [\sFmla{\True}{\formula{A}(a)}, just={\TRule{\True}{\land}[4]}
                [\sFmla{\True}{\formula{B}(a)}, just={\TRule{\True}{\land}[4]}
                ]
              ]
            ]
          ]
        ]
      ]
    ]
  ]
\end{oltableau}
હવે પંક્તિ~$6$ને $\TRule{\True}{\lif}$ લાગુ કરીએ તો !!{tableau} બંધ
થાય છે:
\begin{oltableau}
  [\sFmla{\False}{\lexists[x][\formula{C}(x,b)]}, just=\TAss
    [\sFmla{\True}{\lexists[x][(\formula{A}(x) \land \formula{B}(x))]},
      just=\TAss, checked
      [\sFmla{\True}{\lforall[x][(\formula{B}(x) \lif \formula{C}(x,b))]},
        just=\TAss
        [\sFmla{\True}{\formula{A}(a) \land \formula{B}(a)},
          just={\TRule{\True}{\lexists}[2]}, checked
          [\sFmla{\False}{\formula{C}(a,b)},
            just={\TRule{\False}{\lexists}[1]}
            [\sFmla{\True}{\formula{B}(a) \lif \formula{C}(a,b)},
              just={\TRule{\True}{\lforall}[3]}, checked
              [\sFmla{\True}{\formula{A}(a)},
                just={\TRule{\True}{\land}[4]}
                [\sFmla{\True}{\formula{B}(a)},
                  just={\TRule{\True}{\land}[4]}
                  [\sFmla{\False}{\formula{B}(a)},
                    just={\TRule{\True}{\lif}[6]},close
                  ]
                  [\sFmla{\True}{\formula{C}(a,b)},
                    just={\TRule{\True}{\lif}[6]},close
                  ]
                ]
              ]
            ]
          ]
        ]
      ]
    ]
  ]
\end{oltableau}


\end{ex}

\begin{ex}
આપણે નીચેના ગણ માટે !!a{tableau} રચીએ છીએ:
\[
\sFmla{\True}{\lforall[x][!A(x)]}, \sFmla{\True}{\lforall[x][!A(x)]
  \lif \lexists[y][!B(y)]}, \sFmla{\True}{\lnot\lexists[y][!B(y)]}.
\]
હંમેશની જેમ, ધારણાઓ લખીએ છીએ:
\begin{oltableau}
  [\sFmla{\True}{\lforall[x][\formula{A}(x)]}, just=\TAss
    [\sFmla{\True}{\lforall[x][\formula{A}(x)] \lif
        \lexists[y][\formula{B}(y)]}, just=\TAss
      [\sFmla{\True}{\lnot\lexists[y][\formula{B}(y)]}, just=\TAss
      ]
    ]
  ]
\end{oltableau}
પંક્તિ~$3$ને $\TRule{\True}{\lnot}$ નિયમ લાગુ કરીને શરૂઆત કરીએ છીએ.
``આઇગનચલ-શરતવાળા નિયમો હંમેશાં પહેલાં લાગુ કરો'' એ નિયમનું એક પરિણામ
આ છે: ``આઇગનચલ-શરત વિનાના પરિમાણક નિયમો જરૂરી થાય ત્યાં સુધી મુલતવી
રાખો.'' શાખા પાડતા નિયમોને પણ મુલતવી રાખો.
\begin{oltableau}
  [\sFmla{\True}{\lforall[x][\formula{A}(x)]}, just=\TAss
    [\sFmla{\True}{\lforall[x][\formula{A}(x)] \lif
        \lexists[y][\formula{B}(y)]}, just=\TAss
      [\sFmla{\True}{\lnot\lexists[y][\formula{B}(y)]}, just=\TAss, checked
        [\sFmla{\False}{\lexists[y][\formula{B}(y)]}, just={\TRule{\True}{\lnot}[3]}]
      ]
    ]
  ]
\end{oltableau}
નવી પંક્તિ~$4$ માટે $\TRule{\False}{\lexists}$ જોઈએ, જે આઇગનચલ-શરત
વિનાનો પરિમાણક નિયમ છે. તેથી તેને મુલતવી રાખીને પંક્તિ~$2$ પર
$\TRule{\True}{\lif}$ વાપરીએ છીએ.
\begin{oltableau}
  [\sFmla{\True}{\lforall[x][\formula{A}(x)]}, just=\TAss
    [\sFmla{\True}{\lforall[x][\formula{A}(x)] \lif
        \lexists[y][\formula{B}(y)]}, just=\TAss, checked
      [\sFmla{\True}{\lnot\lexists[y][\formula{B}(y)]}, just=\TAss, checked
        [\sFmla{\False}{\lexists[y][\formula{B}(y)]},
          just={\TRule{\True}{\lnot}[3]},
          [\sFmla{\False}{\lforall[x][\formula{A}(x)]}, just={\TRule{\True}{\lif}[2]}]
          [\sFmla{\True}{\lexists[y][\formula{B}(y)]}, just={\TRule{\True}{\lif}[2]}]
        ]
      ]
    ]
  ]
\end{oltableau}
બંને નવાં !!{signed formula}s માટે આઇગનચલ-શરતવાળા નિયમો જોઈએ છે, તેથી
હવે એ નિયમો લાગુ કરવા જોઈએ:
\begin{oltableau}
  [\sFmla{\True}{\lforall[x][\formula{A}(x)]}, just=\TAss
    [\sFmla{\True}{\lforall[x][\formula{A}(x)] \lif
        \lexists[y][\formula{B}(y)]}, just=\TAss, checked
      [\sFmla{\True}{\lnot\lexists[y][\formula{B}(y)]}, just=\TAss, checked
        [\sFmla{\False}{\lexists[y][\formula{B}(y)]},
          just={\TRule{\True}{\lnot}[3]}
          [\sFmla{\False}{\lforall[x][\formula{A}(x)]}, just={\TRule{\True}{\lif}[2]},checked
            [\sFmla{\False}{\formula{A}(b)}, just={\TRule{\False}{\lforall}[5]}]
          ]
          [\sFmla{\True}{\lexists[y][\formula{B}(y)]}, just={\TRule{\True}{\lif}[2]},checked
            [\sFmla{\True}{\formula{B}(c)}, just={\TRule{\True}{\lexists}[5]}]
          ]
        ]
      ]
    ]
  ]
\end{oltableau}
શાખાઓ બંધ કરવા માટે પંક્તિઓ $1$ અને~$3$નાં !!{signed formula}s વાપરવાં
પડે છે. તેમને અનુરૂપ નિયમો (\TRule{\True}{\lforall} અને
\TRule{\False}{\lexists})ને આઇગનચલ-શરતો નથી, તેથી અનુકૂળ પદો પસંદ
કરવાની આપણને છૂટ છે. આ કિસ્સામાં અનુક્રમે $b$ અને~$c$ અનુકૂળ છે.
\begin{oltableau}
  [\sFmla{\True}{\lforall[x][\formula{A}(x)]}, just=\TAss
    [\sFmla{\True}{\lforall[x][\formula{A}(x)] \lif
        \lexists[y][\formula{B}(y)]}, just=\TAss, checked
      [\sFmla{\True}{\lnot\lexists[y][\formula{B}(y)]}, just=\TAss, checked
        [\sFmla{\False}{\lexists[y][\formula{B}(y)]},
          just={\TRule{\True}{\lnot}[3]}
          [\sFmla{\False}{\lforall[x][\formula{A}(x)]},
            just={\TRule{\True}{\lif}[2]},checked
            [\sFmla{\False}{\formula{A}(b)},
              just={\TRule{\False}{\lforall}[5]}
              [\sFmla{\True}{\formula{A}(b)},
                just={\TRule{\True}{\lforall}[1]},close
              ]
            ]
          ]
          [\sFmla{\True}{\lexists[y][\formula{B}(y)]},
            just={\TRule{\True}{\lif}[2]},checked
            [\sFmla{\True}{\formula{B}(c)},
              just={\TRule{\True}{\lexists}[5]}
              [\sFmla{\False}{\formula{B}(c)},
                just={\TRule{\False}{\lexists}[4]},close
              ]
            ]
          ]
        ]
      ]
    ]
  ]
\end{oltableau}
\end{ex}

\begin{prob}
નીચેના માટે બંધ !!{tableau}s આપો:
\begin{enumerate}
\item $\sFmla{\False}{(\lforall[x][!A(x)] \land \lforall[y][!B(y)])
\lif \lforall[z][(!A(z) \land !B(z))]}$.
\item $\sFmla{\False}{(\lexists[x][!A(x)] \lor \lexists[y][!B(y)])
\lif \lexists[z][(!A(z) \lor !B(z))]}$.
\item $\sFmla{\True}{\lforall[x][(!A(x) \lif !B)]},
\sFmla{\False}{\lexists[y][!A(y)] \lif !B}$.
\item $\sFmla{\True}{\lforall[x][\lnot !A(x)]},
\sFmla{\False}{\lnot\lexists[x][!A(x)]}$.
\item $\sFmla{\False}{\lnot\lexists[x][!A(x)] \lif \lforall[x][\lnot
!A(x)]}$.
\item $\sFmla{\False}{\lnot\lexists[x][\lforall[y][((!A(x,y) \lif
\lnot !A(y,y)) \land (\lnot !A(y,y) \lif !A(x,y)))]]}$.
\end{enumerate}
\end{prob}

\begin{prob}
નીચેના માટે બંધ !!{tableau}s આપો:
\begin{enumerate}
\item $\sFmla{\False}{\lnot\lforall[x][!A(x)] \lif \lexists[x][\lnot!A(x)]}$.
\item $\sFmla{\True}{(\lforall[x][!A(x)] \lif !B)}, \sFmla{\False}{\lexists[y][(!A(y) \lif !B)]}$.
\item $\sFmla{\False}{\lexists[x][(!A(x) \lif \lforall[y][!A(y)])]}$.
\end{enumerate}
\end{prob}

\end{document}
